An observer measured trigonometric parallaxes of stars in a star cluster. Due
to random errors, the measured parallax values are distributed symmetrically
around the expected value with standard deviation equal to 0.05\,mas
(milliarcsec). Assume there are no systematic errors and assume all stars in
the said cluster have the same luminosity. It is known that the distance of
this cluster from us is $R = 5$\,kpc.

He gave the data table to 4 of his students (A, B, C and D) and they estimated
the distance to the cluster in the following ways:

\begin{description}
\item[A.] Convert each parallax measurement into distance and then find the
  average distance ($R_{\mathrm{A}}$)
\item[B.] Take the average of all parallaxes first and then convert the average
  parallax into distance. ($R_{\mathrm{B}}$)
\item[C.] Convert each parallax measurement into distance and then take the
  median distance value. ($R_{\mathrm{C}}$)
\item[D.] Find the median value of the parallaxes and then convert the median
  value into distance. ($R_{\mathrm{D}}$)
\end{description}

State if the following statements are true or false. \textbf{In case a given
mathematical relation is false, give the correct relation.}

\begin{description}
\item[(l)] If the $i^{\mathrm{th}}$ star gave the smallest value of parallax
  and the $j^{\mathrm{th}}$ star gave the highest value of parallax, in all
  likelihood $R_i - R > R - R_j$
\item[(m)] $R_{\mathrm{A}} = R$ (i.e.\ there is a high chance that the distance
  estimated by A fairly matches the true distance)
\item[(n)] $R_{\mathrm{B}} = R$ (i.e.\ there is a high chance that the distance
  estimated by B fairly matches the true distance)
\item[(o)] $R_{\mathrm{C}} < R$ (i.e.\ there is a high chance that the distance
  estimated by C will be systematically lower than the true distance)
\item[(p)] $R_{\mathrm{D}} = R$ (i.e.\ there is a high chance that the distance
  estimated by D fairly matches the true distance)
\end{description}
